% Einheit 2 von 5 – Tangente als Berührung (bank/tangente-normale-schnittwinkel/e2.jsonl, zusammenbau v0.1)
%% TODO Zeitmarke Einheit 2: Marken-Zeile nicht lesbar
%% TODO Zweigzeile Teil 1: was hier gelernt wird (Satz in Schülersprache aus dem Einheitstitel) – Einheit 2
\einheitenkopf[e2]{Einheit 2 von 5 · Tangente als Berührung}
\zweigzeile{Abitur GK}
\setcounter{aufgabe}{12}

%% TODO Titel: Ich-kann-Satz fehlt in der Bank (Platzhalter: Kettenname) – Nr. 13
%% TODO Anweisung: ein Satz über der ersten Teilaufgabe fehlt in der Bank – Nr. 13
\begin{aufgabe}{Berührung nachweisen}
%% TODO \rechenplatz{n}: Schrittzahl des Grundfalls fehlt in der Bank (nur bei fester Schreibform, 2.3 b)
\begin{teile}
\teil An der Stelle $x = 1$ gilt $f(1) = g(1)$, und die Gerade $g$ hat die Steigung $f'(1)$. Wie verhält sich $g$ zum Graphen von $f$ an dieser Stelle? \\ \kreuz{berührt} \\ \kreuz{schneidet nur} \\ \kreuz{weder noch}
\teil Ist $g(x) = -x + 1$ Tangente an den Graphen von $f(x) = x^3 - 2x^2 + 1$ an der Stelle $x = 1$? ja / nein: \leerfeld
\teil $f(x) = x^2 + 1$ und $g(x) = -x^2 + 4x - 1$ – weise nach, dass sich die Graphen in ihrem gemeinsamen Punkt berühren, und gib die gemeinsame Tangente an. $y =$ \leerfeld
\teil Das Profil einer Kübelwand besteht aus zwei Teilen: $f(x) = \frac{1}{4}x^2 + 3$ für $2 \le x \le 4$ und $g(x) = x + 2$ für $-2 \le x \le 2$; bei $x = 2$ stoßen sie aneinander. Weise nach, dass die Wand dort keinen Knick hat. \hfill (FHR 2019)
\teil Die Tangente an den Graphen von $f(x) = x^3 - 9x^2 + 24x - 16$ im Punkt $S(0 | f(0))$ hat einen weiteren gemeinsamen Punkt mit dem Graphen. Koordinaten? \hfill (Abitur 2023 GK) \\ \leerfeld
\teil Die Abbildung zeigt den Graphen von $f$ und den Punkt $Q(0 | -1)$. Lege mit dem Lineal eine Tangente an den Graphen, die durch $Q$ geht, zeichne sie ein und gib ihre Gleichung an. \hfill (Abitur 2020 GK) \\

\begin{ksys}[xmin=-4,xmax=4,ymin=-2,ymax=6] \funktion{0.5*\x^2+1}{f} \punkt{0}{-1}{Q} \end{ksys}
$y =$ \leerfeld
\teil Sara löst das Gleichungssystem $f(u) = 3u - 4$ und $f'(u) = 3$. Welche Frage beantwortet sie damit?
\teil Der Graph von $f(x) = x \cdot e^{2 - x}$ wird nahe $x = 2$ durch die Tangente $t$ an dieser Stelle ersetzt. Gleichung von $t$, und weicht $t(3)$ um weniger als $10\,\%$ von $f(3)$ ab? \hfill (Abitur 2018 LK) \\ $t(x) =$ \leerfeld , ja / nein: \leerfeld
\teil $f(x) = x^2 - 4x + 5$ – weise nach, dass sich die Tangenten an den Stellen $1$ und $3$ im Punkt $S(2 | 0)$ schneiden.
\teil Ein Hügelprofil wird für $-4 \le x \le 4$ durch $f(x) = -\frac{1}{4}x^2 + 4$ beschrieben ($1$ LE = $10$ m). Eine Kamera steht im Punkt $K(5 | 10)$. Ist der Hang von der Kamera aus vollständig einsehbar? Untersuche rechnerisch. \hfill (Abitur 2025 LK) \\ \leerfeld
\end{teile}
\end{aufgabe}

%% TODO Titel: Ich-kann-Satz fehlt in der Bank (Platzhalter: Kettenname) – Nr. 14
%% TODO Anweisung: ein Satz über der ersten Teilaufgabe fehlt in der Bank – Nr. 14
\begin{aufgabe}{Berührpunkt einer Tangente durch einen vorgegebenen Punkt berechnen}
\begin{teile}
\teil Ein Glas hat den Längsschnitt $f(x) = -\frac{1}{16}x^4 + x^2$ für $-4 \le x \le 4$ ($1$ LE = $1$ cm). Ein gerader Halm steht mit dem unteren Ende im Punkt $P(-1 | f(-1))$ und berührt das Glas im Punkt $Q(u | f(u))$ mit $u > 0$. Berechne $u$ auf zwei Stellen. \hfill (Abitur 2017 LK) \\ $u \approx$ \leerfeld
\end{teile}
\end{aufgabe}

%% TODO Titel: Ich-kann-Satz fehlt in der Bank (Platzhalter: Kettenname) – Nr. 15
%% TODO Anweisung: ein Satz über der ersten Teilaufgabe fehlt in der Bank – Nr. 15
\begin{aufgabe}{Zeitpunkt bei gleichbleibender Änderungsrate über den Schnitt der Tangente mit der Zeitachse grafisch bestimmen}
\begin{teile}
\teil Die Abbildung zeigt das Wasservolumen $V$ in einem Becken ($x$ in Stunden, $V(x)$ in m³). Die Änderungsrate zum Zeitpunkt $x = 10$ bleibt ab dann erhalten, bis das Becken leer ist. Beschreibe, wie man diesen Zeitpunkt grafisch bestimmt, und bestimme ihn. \hfill (Abitur 2017 LK) \\

\begin{ksys}[xmin=0,xmax=30,ymin=0,ymax=45,xstep=5,ystep=5,xlabel=x in h,ylabel=V in m³,ablesen] \funktionab{40-0.1*\x^2}{V}{0}{20} \end{ksys}
nach \leerfeld[h]
\end{teile}
\end{aufgabe}

%% TODO Titel: Ich-kann-Satz fehlt in der Bank (Platzhalter: Kettenname) – Nr. 16
%% TODO Anweisung: ein Satz über der ersten Teilaufgabe fehlt in der Bank – Nr. 16
\begin{aufgabe}{Lage zweier Graphen mit gemeinsamen Endpunkten über die Steigungen in den Endpunkten begründen}
\begin{teile}
\teil Zwei Seile hängen zwischen den Punkten $(-4 | 3)$ und $(4 | 3)$, das eine entlang $g(x) = \frac{1}{8}x^2 + 1$, das andere entlang $h(x) = \frac{1}{128}x^4 - \frac{1}{16}x^2 + 2$. Berechne die Steigungen beider Seile in den Befestigungspunkten und begründe damit, dass $h$ weder ganz oberhalb noch ganz unterhalb von $g$ verläuft. \hfill (Abitur 2018 LK)
\end{teile}
\end{aufgabe}

%% TODO Titel: Ich-kann-Satz fehlt in der Bank (Platzhalter: Kettenname) – Nr. 17
%% TODO Anweisung: ein Satz über der ersten Teilaufgabe fehlt in der Bank – Nr. 17
\begin{aufgabe}{Berührung nachweisen – Fehler finden, Begründen}
\begin{teile}
\teil Ich finde den Fehler: Ist $g(x) = 2x + 1$ Tangente an $f(x) = x^2 - 2x + 3$ an der Stelle $2$? Jonas rechnet: \rechnung{f'(x) &= 2x - 2 \\ f'(2) &= 2 = \text{Steigung von } g \quad \Rightarrow \text{ Tangente}}
\teil Begründe, warum die Berührstelle eine doppelte Nullstelle von $f(x) - t(x)$ ist.
\end{teile}
\end{aufgabe}
